\documentclass[book.tex]{subfiles}
\begin{document}

\chapter{Differentiable Fluid Dynamics Solvers}
\label{ch:fluids_modeling}

\label{chap:ml_coarse_fluid}
\noindent\rule{11cm}{0.4pt} \\
\textbf{Prerequisites:} \autoref{chap:qft_basics}  \\
\textbf{Difficulty Level:} ** \\
\noindent\rule{11cm}{0.4pt}
\vspace{5mm}

The Navier-Stokes equations can be directly differentiated through to facilitate approximate solution of fluid systems. These methods can be applied for both direct numerical simulation approaches and large eddy simulation approaches. The central idea is to use a much coarser grid than is usually admissible for simulations, but to use a machine-learned component to adapt for the loss of resolution\cite{kochkov2021machine}.

\section{Learned Solver}
The core idea is to implement a differentiable finite volume solver. The full model is trained simply with $L^2$ loss between predictions and ground truth data obtained from DNS solution. 

\section{Evaluations}

The energy spectrum
\begin{align}
E(k) &= \frac{1}{2}|u(k)|^2
\end{align}
can be computed and compared to the energy spectrum of DNS data. Models can also be tested on conditions not originally considered. The scaling of this coarse grained approach is roughly given by
\begin{align}
T \sim (C_{ML} + C_{Physics})\left ( \frac{N}{k} \right )^{d+1}
\end{align}
where $N$ is the size of the discretized grid, and $C_{ML}$ and $C_{Physics}$ are the costs of the machined-learned and physical updates. The ML updates are for now much slower than the physical updates.

%"Capturing the essential dynamics of complex physical systems such as climate, weather or turbulent flows with a small number of computational degrees of freedom is a long-standing problem in applied mathematics. Recently, there has been a flurry of interest in using machine learning to accomplish this. In this talk I'll discuss the approach and perspectives that have been driving our team in this line of research. I will present the results from our recent paper on accelerating turbulent flows in two dimensions, where we develop a simulation method that combines ML and CFD techniques to accurately operate on coarser grids. For both direct numerical simulation and large eddy simulation, our method is as accurate as baseline solvers with 8-10x finer resolution in each spatial dimension, resulting in 40-80x fold computational speedups. Accompanying the paper, we have open-sourced jax-cfd - a differentiable fluid dynamics solver which I'll discuss during the code walkthrough section."


\end{document}